Mario Leoncini

ELEMENTI

DI STRATEGIA

NEGLI SCACCHI

\_\_\_\_\_\_\_\_\_\_\_\_\_\_\_\_\_\_\_\_\_\_\_\_\_\_\_\_

1992. Prima edizione su carta 2004

\textbf{INTRODUZIONE}

Se la tattica è lo sfruttamento combinativo di una debolezza e la
tecnica è la capacità di vincere una partita considerata vantaggiosa, la
strategia è il piano di gioco.

Qualcuno, in modo divertente quanto efficace, ha sintetizzato la
differenza tra strategia e tattica dicendo: \emph{``La tattica è sapere
che cosa fare quando c'è qualcosa da fare e la strategia è sapere che
cosa fare quando non c'è niente da fare''.}

La buona accoglienza riservata a questo scritto, nato alla fine del 1992
in forma di lezioni per ``Mc link'', il BBS della nota rivista di
informatica ``McMicrocomputer'', mi ha indotto a raccogliere il
materiale e ampliarlo in un libro. Libro che si occupa di strategia e
più in particolare degli elementi rilevanti per la formulazione del
piano di gioco. La strategia ha per scopo la formazione di debolezze
tali da poter essere sfruttate con colpi tattici o in sede tecnica.

La sua formulazione deve naturalmente tener conto di tutti gli elementi
presenti sulla scacchiera che però, beninteso, varia al variare della
posizione. Non è quindi possibile formulare un piano unico valido per
tutte le stagioni ma è anche vero che, data una posizione, talvolta
possono formularsi più piani. La scelta del piano dipende allora
dall'indole del giocatore ma non si possono formulare piani che non
tengano conto degli elementi strategici presenti nella posizione.

Un altro errore comune è credere che la differenza tra un maestro e un
principiante risieda nella capacità di calcolo; certo, un bravo
giocatore è capace di calcolare con precisione anche lunghe varianti; ma
la superiorità è soprattutto di ordine strategico. Il maestro \emph{sa}
che cosa fare in qualsiasi posizione senza bisogno di calcoli
approfonditi. La formulazione di un piano riduce drasticamente la
necessità del calcolo delle varianti; in questo senso nel corso della
partita il maestro può calcolare meno di un principiante ma vincere lo
stesso. È per questo che i forti giocatori possono giocare contro molti
avversari contemporaneamente in simultanea e batterli.

In verità, chiunque giochi a scacchi, fosse anche alle primissime armi,
dà al proprio gioco almeno un simulacro di strategia. Nella versione più
semplice essa consiste nel domandarsi, dopo ogni mossa dell'avversario:

1) se il proprio materiale è minacciato e si è sotto minaccia di scacco
matto; se la risposta a questa domanda è sì, allora

1,1) porvi rimedio; altrimenti

2) vedere se si può, a nostra volta, minacciare materiale o il matto.

Per ragionare in questi termini occorre naturalmente conoscere il valore
dei pezzi.

Se si prende come unità di misura il pedone, il loro valore è, grosso
modo, questo: Donna~=~9, Torre = 5, Alfiere = 3, Cavallo = 3.

Ecco dunque che lo studio degli elementi strategici è anche lo studio
dei fattori che vanno considerati, oltre al materiale, per capire chi è
in vantaggio data una qualsiasi posizione. Di più: non solo il materiale
non è l'unico elemento, come può esserlo in giochi più semplici degli
scacchi, ma il suo valore è "relativo", varia a seconda della
disposizione degli altri pezzi.

Questo libro può essere letto da chi già conosce le regole e ha un
minimo di esperienza ma si rivolge, in particolare, a quel vasto settore
di giocatori che desiderano migliorare il proprio gioco e che sono
collocati tra la terza categoria sociale e il titolo di candidato
maestro. Costoro sono soliti commettere l'errore di tralasciare lo
studio della strategia in favore di quello delle aperture. Non nego che
lo studio delle aperture (e dei finali, anch'essi spesso immeritatamente
trascurati) sia indispensabile, ma rischia di trasformarsi in inutile
esercizio mnemonico se non si capiscono i motivi che si nascondono
dietro le successioni di mosse, presentate, nei monumentali libri tipo
``Enciclopedia'' di Belgrado, in forma inevitabilmente asettica.

Ci sono giocatori che entrano in profonda crisi, e che commettono gli
errori più incredibili, quando, nella variante d'apertura preferita, il
loro avversario esegue una mossa non studiata o non contemplata dai
``sacri'' testi. Questi giocatori troveranno qui raccolti, e spiegati in
modo sistematico, elementi che hanno cominciato a intravedere e intuire
fin dalle prime partite.

La loro conoscenza e il loro studio sono gli indispensabili ingredienti
per un gioco di livello magistrale.

Mario Leoncini 1992--1996
